%
%      Brno University of Technology
%      Faculty of Information Technology
%
%      MSc Thesis	2009/2010
%
%      Design of S-Boxes Using Genetic Algorithms
%
%      Bedrich Hovorka
% -----------------------------------------------------
%
% Design description and various options - enough to mention

\section{Summary of Technical and Functional Requirements}
The objective of this work is to implement a program that searches for S-boxes using evolutionary algorithms,
and subsequently perform experiments with this program.
The implemented program has the following main technical requirements, listed
in descending order of importance:
\begin{description}
  \item[Easy experiment definition] When describing experiments, it is necessary
  to avoid routine steps, such as complicated declarations and time-consuming code compilation, ideally
  making it possible to specify and evaluate experiments interactively.
  \item[Speed] The execution of evolution and calculation of criteria must be performed as quickly as possible.
  \item[Portability] The program should be able to be compiled and run on the widest possible range of platforms used
     for scientific computing on standard personal computers.
\end{description}

\subsection*{Functional requirements:}
\begin{description}
  \item[Various S-box representations] To find S-boxes in the most suitable representation (see subsection \ref{secSboxTypes}).

  \item[Various types of search algorithms] For the purposes of comparing the applicability of EA (see chapter \ref{evo})
  it is useful to compare them with random search.

  \item[Search criteria] The required search criteria (see subsection \ref{secCriterions})
  are parameters of the search algorithms.

  \item[Experiment execution] The actual execution of experiments should have a unified application interface.
  The ability to first define multiple experiments and then run them all at once (e.g., in parallel).

  \item[Statistics collection] To evaluate individual runs, it is necessary for the program to record statistics during evolution.
  Evolutionary experiments are often evaluated using so-called \uv{box plots}, which require calculating
  the minimum, maximum values and quartiles, fitness deviation at a fixed number of generations.
  Or \cite{millerThomson} by comparing what effort (e.g., number of generations, population size, runtime) with what quality of results (highest fitness)
  was expended.
\end{description}

\section{Prototype Application Design}
After establishing the main requirements, an application prototype was first assembled
to clarify how much needs to be implemented,
in a form similar to UML diagrams; it contains classes with empty methods.
The superclass of evolutionary algorithms maintains a population of candidate solutions. It contains the general logic for
the evolution process and statistics collection. Figure \ref{evolutionClassDiagram} illustrates the main evolutionary
algorithms.
\insertimage{evolutionClassDiagram}{Prototype class diagram of evolutionary algorithms}

Each search algorithm has its specific population modification in each generation, defined
by the method \sourceRef{variation()}.
During object construction (before running \sourceRef{evolution()}), the S-box representation type and evolution
parameters are passed. The sets of criterion evaluators are also passed during construction.
For the EDA algorithm, the important attribute is \sourceRef{model()}, which represents the statistical model
(subtype of EDA) that will be used.
The simplicity of the diagram proved important in the implementation phase, when based on it,
an interface between Python and C++ was created and further elaborated.

\section{Various Types of S-box Representation}\label{secSboxTypes}
Some evolutionary algorithms may require a specific representation of the substitution box.
Some of the representations may also allow optimal implementation if they are models of the target system.
Each of them has its specific initialization, mutation, and crossover.

\subsection{Output List}
Most commonly, S-boxes are represented as an output list (array) of the size of all
input combinations, where the S-box input is the index and the output is the element
at that position. If the S-box is to be bijective, then this list represents a permutation.
For example
\begin{center}
2 0 3 1
\end{center}
is a $2 \times 2$ S-box that is bijective.
With this representation, it is possible to choose whether bijectivity should be required.
The genetic operators are determined accordingly. For bijective, mutation is performed
by swapping genes. Crossover proceeds as in permutation problems based on order.
If bijectivity were not required, single-point crossover would be used, as in the following
representation, but with a different set of alleles.

\subsection{Binary Representation}
Binary representation consists of Boolean functions for each input.
The bit of the Boolean function at the $i$-th position represents the output when the input is $i$.
The S-box from the previous representation can be written in binary as follows (columns still represent the same value):
\begin{center}
0 0 1 1\\
1 0 1 0
\end{center}
And each row represents a Boolean function.
Incidentally, this S-box can also be represented as polynomials of individual functions,
where the variables are input bits:
$$s(x)_0 = b , s(x)_1 = 1 \oplus a$$
For genetic operators, the list of Boolean functions is treated as a single chromosome, on which
single-point crossover is performed and mutation changes a bit at a random position.

\subsection{Acyclic Gate Graph}
This representation comes from Cartesian Genetic Programming (notation see appendix \ref{cgpfunctions})
and is suitable for finding optimal hardware implementation.
The S-box from the previous representation can be described, for example, by a chromosome in the form
\begin{verbatim}
{2,2,2,2,2,1,1}([2]1,0,5)([3]1,0,4)([4]1,3,3)([5]3,3,1)(1,2)
\end{verbatim}
which is also shown in Figure \ref{cgpExample}.
The first (index 0) output bit is connected to the second input. And the second output bit is the negation
of the first input, which corresponds to the polynomials $s(x)_0 = b , s(x)_1 = 1 \oplus a$.
\insertimage[0.6\textwidth]{cgpExample}{Example of acyclic gate graph}
In this representation, only mutation satisfying the conditions already mentioned
in subsection \ref{algorithmCGP} is permitted.

\subsection{Sequence of Triple Instructions for i686 Platform}\label{secAbouti686}
This representation was inspired by the implementation of the SubCrumb algorithm Luffa \cite{specLuffa}.
With the modification that it is designed for bits at inputs and outputs.
The number of input and output bits are multiples of four.
The only structural parameter is the number of bit quartets $p$ at inputs and outputs (the number of inputs and outputs is equal).
The representation is a model of i686 series processors, for the purpose of finding optimal implementation on
these processors -- optimal in software. The model consists of five registers, each storing $p$ bits.
The computation of each S-box consists of six cycles. In each cycle there are three instructions, which
are assumed to be executed by the processor in parallel, therefore the same
register must not repeat at the instruction output (CREW parallelism). The instructions that can occur in each cycle are
XOR, AND, OR, NOT. Additionally, in the first cycle it is possible to execute MOV and in the last cycle NOP.

The computation process begins by dividing the input into four parts $a_0$, $a_1$, $a_2$, $a_3$.
Each part $a_i$ has $p$ bits and is loaded into register $r_i$, $r_4 = 0$.
In the simulation of one cycle, the register values are first stored in temporary variables and by executing individual
instructions, new values are set. The S-box output consists of $p$-bit parts $y_i$
as follows: $y_0 = r_0$, $y_1 = r_1$, $y_2 = r_3$ and $y_3 = r_4$.

For the output list
\begin{center}
2 9 13 13 2 9 13 13 2 9 13 13 13 6 6 6
\end{center}
the sequence of triple instructions is
\begin{center}
\begin{verbatim}
{AND r2, r3 ; MOV r3, r1 ; OR r4, r0}
{XOR r1, r3 ; XOR r4, r4 ; OR r0, r1}
{AND r1, r3 ; NOT r4 ; AND r3, r3}
{OR r2, r3 ; XOR r0, r2 ; AND r4, r1}
{OR r2, r3 ; OR r0, r4 ; OR r3, r0}
{XOR r0, r1 ; NOT r1 ; OR r4, r0}
\end{verbatim}
\end{center}

And for the output list
\begin{center}
15 8 15 8 15 3 15 3 15 8 1 3 15 3 1 8
\end{center}
the sequence of triple instructions is
\begin{center}
\begin{verbatim}
{AND r2, r2 ; AND r3, r1 ; OR r1, r2}
{AND r3, r4 ; AND r2, r0 ; XOR r4, r3}
{XOR r4, r2 ; NOT r0 ; NOT r2}
{AND r1, r4 ; NOT r3 ; AND r0, r2}
{XOR r1, r2 ; NOT r4 ; XOR r0, r1}
{NOP ; OR r3, r3 ; AND r2, r4}
\end{verbatim}
\end{center}
Instructions are in the form \verb|INST output, input|.
For NOT, the input is also the output; the first operand is used for its definition.
And NOP has no operands. Thus unnecessary operands are ignored.

Crossover occurs at the cycle level. Single-point crossover is performed with cycles from parents.
During mutation, all cycles are traversed. In each cycle, with probability $pMut$ an
instruction is selected for modification. If an instruction is selected, one of three possible changes is chosen
with uniform probability distribution:
\begin{enumerate}
  \item change instruction type -- one of the allowed instructions
  in the current cycle is selected with uniform probability distribution
  \item change second operand -- one of five possible registers is selected with uniform distribution
  \item change first operand -- similarly, but here a register must not be selected that is already in the current cycle
  used as the first operand in another instruction
\end{enumerate}

Table \ref{tabReprezantationUsability} shows which search algorithm can use which representation from the perspective of the strict definition
of these algorithms.
For acyclic gate graphs, crossover is not defined, so it must be disabled.
If an EDA algorithm works with numeric alleles, it can also search for a numeric list, however
that statistical model will be unnecessarily complicated. VEGA and SPEA algorithms actually handle what GP does;
this table takes into account the need for multi-criteria optimization.
\begin{table}[t]
  \begin{center}
     \begin{tabular}{| l | c | c | c | c |}
     \hline
     Algorithm & Binary & Output list & CGP & i686 \\ \hline
     basic GA & suitable & suitable & impossible & impossible \\ \hline
     GP & possible & possible & possible & suitable \\ \hline
     EDA & suitable & unsuitable & impossible & impossible \\ \hline
     Random & possible & possible & possible & possible \\ \hline
     ParallelRandom & possible & possible & suitable & possible \\ \hline
     VEGA & suitable & suitable & suitable & suitable \\ \hline
     SPEA & suitable & suitable & suitable & suitable \\ \hline
     \end{tabular}
     \caption{Usability of S-box representations in individual search algorithms}
     \label{tabReprezantationUsability}
  \end{center}
\end{table}

\section{Languages and Platform Used}
To meet the basic technical requirements, a combination of Python and C++ was chosen.


\subsubsection{Python}
The advantage of Python is the quantity of libraries written in it, and the ease of their cooperation.
These include, for example, pickle for simple persistent storage of almost any object, which
can be used for continuous saving of partial experiment results.
The numpy library is primarily for working with large vectors and matrices with many mathematical functions.
These functions are written in C, and it can easily work with this language.
In addition, there is the ctypes library, which allows working with C data types
in functions that are part of a dynamically linked library to Python.
For writing experiments, the advantage is the interactivity of command entry, the so-called \uv{ipython}
interpreter, which in combination with numpy and a graph plotting library is an alternative to
Matlab.
The disadvantage is slow code execution, which must be addressed by implementing time-critical functions in C or
C++.

\subsubsection{C++}
C++ can generate fast code for many processors and platforms. Written code is
portable when following certain rules. It is suitable for implementing computations that
need to run quickly. Many libraries with solved algorithms also exist for it.
The disadvantage is complicated declarations and definitions, which slow down experiment descriptions. % to musi uznat i PP

\pagebreak

\subsubsection{Galib}
Galib is a library originally for genetic algorithms, but with its general design it can handle more general classes,
for example genetic programming. It is designed for C++. Other evolutionary algorithms must be implemented additionally.
The Galib library best met the previous requirements. Most important was generality for comparison purposes.
This library \cite{galib} allows defining custom search algorithms, custom operators, or
custom genome representation, etc. During programming, work is done with two main classes
representing the genome and the second representing the genetic algorithm.
Each genome instance represents one problem solution.
The genetic algorithm uses an objective function (user-defined)
to determine the genome survival criterion (fitness). Common genetic operators (selection, mutation, crossover)
are built-in. During evolution, it collects statistics. Evaluators and operators can be changed during runtime
of evolution.
On the other hand, its library functions work correctly only during sequential execution,
if parallelization is needed, its code must be modified. Many methods are not virtual
and when using custom components, it is additionally necessary
to perform type casting\footnote{They even advise this in the examples}.

\subsubsection{OpenMP}
OpenMP is a specification for parallelizing algorithms using directives that
compilers translate according to the target platform and processor.
OpenMP is designed for various languages; for example, in C++ directives are specified using
\verb|#pragma omp|.
GCC compilers implement
some of its directives for C/C++ languages using POSIX threads. These directives are designed so that on single-processor
machines, source code can be compiled without the switch enabling OpenMP, in such a way that
these directives will simply be ignored. This is one of the portability features
that this specification has. These directives find their greatest application primarily
in parallelizing existing sequential code. OpenMP automatically attempts to
utilize all available processors or cores for parallel execution of code marked
with directives in a certain way (parallel loop, task), so they could be executed in parallel.
The programmer often must use these directives to mark
shared and local variables. Other directives such as barrier and critical section allow
performing synchronization between threads.

\subsubsection{Discussion of Other Program Design Variants}
As indicated above, the entire project could be implemented only in Python using the mathematical system
sagemath, which already contains functions for working with S-boxes and calculations of their criteria.
However, the implementation of evolutionary algorithms would not be efficient due to the high-level nature of the language.
%Nebo celé v C++ ale to bych si odstrelil obe nohy,
% ale furt mensi zlo nez Squeak Smalltalk, tam jsem se strelil do hlavy
Demanding computations could also be performed on GPU.
First, however, it is good to have a CPU implementation for comparison, whether GPU implementation would be worthwhile at all.
